%% Curriculum Vitae -- Giosù Aiello -- versione italiana
%% Compilare con pdfLaTeX (lo stile è in cvstyle.sty, la foto in photo.jpg)
\documentclass[11pt,a4paper]{article}
\usepackage[italian]{babel}
\usepackage{cvstyle}

\hypersetup{pdftitle={Curriculum Vitae - Giosù Aiello},
            pdfauthor={Giosù Aiello}}

\begin{document}
\begin{cvlayout}

%% ====================== COLONNA SINISTRA (barra laterale) ==================
\cvname{Giosù Aiello}
\cvphoto{photo.jpg}

\sidebarsection{Competenze Professionali}
\begin{cvskills}
  \item Lavoro di squadra (Teamwork)
  \item Comunicazione efficace
  \item Progettazione e presentazione di progetti
  \item Scrittura scientifica e tecnica
  \item Analisi dei dati con Python e MATLAB
  \item Orientamento alla soddisfazione degli stakeholder
  \item Programmazione: Python, JavaScript, C++, \LaTeX
  \item E altro\dots
\end{cvskills}

\sidebarsection{Lingue}
\textbf{Italiano} (Madrelingua)\\
\textbf{Inglese} (B2)

\sidebarsection{Informazioni Personali}
Cittadinanza: \textbf{Italiana}\\[2.9bp]
Residenza: \textbf{Pisa, Italia}

\sidebarsection{Contatti}
{\small
\MVAt{} \href{mailto:giosue.tanz@gmail.com}{giosue.tanz@gmail.com}\\[6.3bp]
\textbf{Sito Web:}\\
\href{https://giosueaiello.com}{giosueaiello.com}\\[6.3bp]
\textbf{GitHub:}\\
\href{https://github.com/Giosue-tanz}{github.com/Giosue-tanz}
}

%% ======================= COLONNA DESTRA (contenuto) ========================
\cvmaincolumn

\cvsection{Profilo}
\begin{cvprose}
Studente universitario in Fisica presso l'Università di Pisa, con una solida
base accademica e un costante impegno verso l'eccellenza analitica. Il mio
percorso accademico è guidato da un profondo interesse per la comprensione
delle leggi fondamentali della natura e per la loro successiva applicazione
nelle tecnologie all'avanguardia.

Ho sviluppato un solido background sia nella fisica teorica che in quella
sperimentale, con un focus specifico sui sistemi quantistici e sulla
modellazione computazionale. Il mio approccio integra il rigore matematico con
una risoluzione proattiva dei problemi, mirata a contribuire alla ricerca
innovativa.

Oltre al piano di studi ordinario, sono attivamente impegnato in progetti di
ricerca che coinvolgono reti neurali profonde (deep neural networks).
\end{cvprose}

\cvsection{Esperienza}

\cvevent[\cvurl{https://giosue-tanz.github.io/smart-light.it/}{SmartLight}]%
        {CEO \& CTO}{2025--Settembre 2026}
\cvitem{Alla guida dell'iniziativa SmartLight, un progetto basato sull'IA
focalizzato sull'ottimizzazione del flusso del traffico urbano a Pisa. Il
sistema di gestione adattivo sfrutta il Reinforcement Learning e la computer
vision per controllare i segnali stradali in tempo reale, riducendo la
congestione e il consumo energetico. Progetto \textbf{vincitore del contest di
idee per la mobilità sostenibile} promosso dal
\cvurl{https://www.unipi.it/news/due-idee-per-una-mobilita-sostenibile-e-inclusiva-premiati-i-vincitori-del-contest-promosso-dal-contamination-lab/}{Contamination
Lab dell'Università di Pisa}.}

\cvevent{Autore di Appunti di Studio Universitari}{2024--Presente}
\cvitem{Autore e curatore di appunti di studio universitari completi e risorse
accademiche per i corsi fondamentali di Fisica e Matematica, redatti in LaTeX
per facilitare la padronanza concettuale.}

\cvsection{Formazione}

\cvdegree{Studente Universitario in Fisica (Laurea Triennale)}%
         {Università di Pisa}{2023--Presente}

\cvdegree{Diploma di Liceo Scientifico (Opzione Scienze Applicate)}%
         {Liceo Scientifico Enrico Boggio Lera, Catania}{2017--2022}

\end{cvlayout}
\end{document}
